\section{Introduction}\label{sec:introduction}

Operating-system policies have traditionally been designed for generality. They are expected to work reasonably well across many workloads and hardware platforms, and therefore expose only coarse tuning knobs, such as enabling hugepages or adjusting task priorities. However, such a design is increasingly mismatched to modern systems, where applications evolve rapidly, workloads are highly dynamic, and a single machine may combine heterogeneous cores, memory tiers, and accelerators.

As a result, OS policy design is becoming increasingly specialized. For example, recent work customizes OS subsystems via eBPF~\cite{kflex-sosp24,cache-ext-sosp25,sched-ext,meta-lavd} and develops learning-driven OS policies~\cite{heimdall-eurosys25,orca-sigcomm20,kml-tos23,kleio-hpdc19}. This trend is likely to accelerate as ML-based tuning of OS policy knobs matures~\cite{tuning-pacmi25,agentic-tuning-ml4sys25} and, perhaps more crucially, AI agents are increasingly used to generate new kernel policies~\cite{policysmith-hotnets25,glia-arxiv25,barbarians25,vulcan-arxiv25,agentic-os-ml4sys25}. In other words, the OS is shifting from broadly tuned default policies toward mechanisms that are explicitly tailored to particular workloads, hardware characteristics, and deployment assumptions.

Unfortunately, such specialization creates a new class of problems. These policies derive their performance from assumptions about workloads and environments, yet those assumptions are usually neither enforced nor checked at deployment time. A policy may assume that profiled access patterns remain stable, that a learned predictor continues to see in-distribution inputs, or that a particular assumed bottleneck remains dominant. \aditya{forward pointer to examples} When such assumptions fail, the policy may still execute correctly, but it may silently stop delivering the performance benefit that motivated it. End-to-end benchmarking may reveal that a policy underperformed, but it rarely explains why. In particular, it does not identify which assumption failed, under what execution conditions the failure arose, or, crucially, what part of the policy needs to change and how. As a result, policy development often degenerates into a painful trial-and-error cycle.

In this paper, we argue that desired performance properties should be treated as a first-class system-design concern, rather than as something to be evaluated only after a policy has been built. To this end, we propose a framework for \emph{performance introspection}, which equips the system with the ability to check whether an OS policy is actually delivering its intended performance at deployment time and respond when it is not.

Concretely, we realize this idea through the concept of \emph{OS guardrails}. A guardrail pairs a runtime-checkable performance condition with a corrective action that is invoked when the condition is violated. Rather than relying solely on offline benchmarking, guardrails let the OS continuously check whether \aditya{make sure this statement is correct. are guardrails only about assumptions} the assumptions underlying a policy still hold during execution. When they do not, guardrails can identify the relevant problematic behavior and trigger a targeted response, such as updating a parameter, switching policies, or falling back to a safer default.

However, realizing this idea as a practical OS mechanism is not straightforward and raises four core challenges:
\begin{itemize}[leftmargin=*]
    \item \textbf{C1: Performance is hard to observe directly.} The performance objective that a designer cares about is often not directly monitorable inside the kernel. Guardrails therefore need runtime conditions that are observable, yet still faithfully reflect the intended performance goal.
    
    \item \textbf{C2: No single guardrail fits all policies.} Different policies depend on different assumptions and can go wrong in different ways. A memory-management policy, a load balancer, and a learned scheduler may require different \aditya{better word than signal?} signals, different checks, and different repairs.
    
    \item \textbf{C3: Local checks do not automatically imply global correctness.} A runtime check is only a proxy for the end-to-end property the designer ultimately wants. The system must therefore connect local conditions and local repairs to higher-level performance guarantees.
    
    \item \textbf{C4: Introspection must not become the bottleneck.} Since guardrails execute in the OS itself, they must remain lightweight enough to run continuously without undermining the performance they are meant to protect.
\end{itemize}

We present a framework for practical performance introspection for OS policies. We develop principles for designing guardrails whose conditions are monitorable at runtime yet still capture policy-relevant performance behavior. We then introduce a Guardrail Description Language (\lang) for specifying policy-specific checks and corrective actions over kernel-visible signals. To connect these local mechanisms to the higher-level objectives that designers actually care about, we provide a lightweight verification interface that relates guardrail conditions to performance properties and supports reasoning about multiple guardrails simultaneously. Finally, we realize this framework through an eBPF-based implementation that enables low-overhead monitoring in the kernel.

We evaluate guardrails across both traditional kernel policies and learned policies. Our results show that guardrails can (a) detect when deployment-time assumptions break, (b) identify which signals or decisions indicate that a policy is no longer operating in the regime for which it was designed, and (c) enable targeted corrections or safer fallbacks. More broadly, our case studies suggest that performance introspection provides a systematic way to design, debug, and deploy specialized OS policies whose assumptions are checked continuously during execution rather than only after performance has already regressed.